% Contenido del póster: Fasano, Nania y Staglianò (2026), Journal of Economics and Business. DOI 10.1016/j.jeconbus.2026.106343
% Elaborado con apoyo de IA a partir del resumen oficial del artículo (OpenAlex), 2026-11-12.
\def\DOI{10.1016/j.jeconbus.2026.106343}
\def\TituloEs{Predicción de retrasos de pago en el crowdfunding inmobiliario italiano: evidencia con aprendizaje automático}
\def\TituloOriginal{Predicting repayment delays in Italian real estate crowdfunding: Evidence from a machine learning approach}
\def\Autores{Fasano, Nania y Staglianò}
\def\Anio{2026}
\def\Revista{Journal of Economics and Business}

\def\Problema{En el crowdfunding inmobiliario, los retrasos en la devolución del capital son un riesgo central para los inversionistas, y las plataformas necesitan detectarlo a tiempo.}
\def\Objetivo{Desarrollar un modelo de aprendizaje automático para detectar de forma temprana el riesgo de retraso en el pago de los proyectos de crowdfunding inmobiliario.}
\def\Metodo{Datos por proyecto de una plataforma italiana de préstamos. Desde las teorías de la selección adversa y la señalización, las características del prestatario, del proyecto y los indicadores de riesgo de la plataforma se tratan como señales que reducen la asimetría de información. Se comparan cuatro algoritmos de clasificación y se elige \textit{Gradient Boosting}; la interpretación se apoya en valores SHAP.}
\def\Resultados{\begin{itemize}
  \item \textit{Gradient Boosting} se elige como modelo principal por su desempeño predictivo fuera de muestra.
  \item El modelo se complementa con un sistema continuo de puntaje de riesgo para monitorear la cartera.
  \item Los valores SHAP muestran qué variables del proyecto pesan más en el riesgo de retraso.
\end{itemize}}
\def\Conclusion{Evaluar el riesgo crediticio según el momento de pago aporta una herramienta útil para plataformas e inversionistas: más transparencia, gestión anticipada del riesgo y asignación de recursos basada en datos.}
\def\Relevancia{La SMV no publica retrasos ni impagos por proyecto del financiamiento participativo, por lo que en el Perú no puede replicarse este análisis. El estudio muestra el valor de que las plataformas y el supervisor publiquen información de pagos, sobre todo en un mercado que en 2025 financió \cifra{proyFin2025} proyectos.}
\def\PalabrasClave{crowdfunding inmobiliario, retrasos de pago, aprendizaje automático, selección adversa, señalización}
% --- Ficha del estudio (herramientas/posteres/ficha_articulos.py) ---
\def\Ambito{Italia (plataforma de préstamos inmobiliarios)}
\def\Enfoque{Aprendizaje automático (\textit{Gradient Boosting} y SHAP)}
\def\FiguraPeru{1}
% --- fin de la ficha ---
